Reinforcement Learning (RL)~\cite{sutton2018rl} is a framework where an agent
learns by interacting with an environment: at each step, it observes the current
state, chooses an action, receives a reward, and transitions to a new state.
The goal is to maximize the cumulative reward over time.

Deep Q-Learning (DQN)~\cite{mnih2015dqn} solves this by training a neural network
to predict the quality of each possible action in a given state — the so-called
$Q$-value. During training, the agent stores its experiences $(s, a, r, s')$ in
a replay buffer. Periodically, it samples a batch of past experiences and updates
the network to make its $Q$-value predictions closer to the Bellman target:
$r + \gamma \max_{a'} Q(s', a')$, i.e., immediate reward plus the best possible
future value. Two mechanisms stabilize this process: a separate \emph{target
network} that changes slowly, and \emph{experience replay} that breaks
correlations between consecutive samples. Extensions such as Double
DQN~\cite{vanhasselt2016double} reduce overestimation by using separate networks
for action selection and evaluation, while Dueling DQN~\cite{wang2016dueling}
splits the output into a state-value stream and an action-advantage stream for
more efficient learning.

Variational Quantum Circuits (VQCs) are the quantum counterpart of neural
networks: classical data is encoded into a quantum state via rotation gates,
trainable gates transform the state through repeated layers, and measurements
extract classical output features. Unlike classical networks, VQCs suffer from
the barren plateau problem~\cite{mcclean2018barren,cerezo2021cost}, where
gradients vanish exponentially with circuit size, making training difficult.
Data re-uploading~\cite{perez2020reuploading} — repeating the encoding at each
layer — increases expressivity by giving the circuit more access to the input.

Table~\ref{tab:related} summarizes prior work on VQC-based reinforcement
learning. The present work differs by including Double DQN and Dueling DQN
variants with a VQC backend, reporting batched QNode execution as a practical
optimization, and including a supervised distillation experiment to isolate
representational capacity from reward-based learning.

\begin{table}[t]
  \centering
  \caption{Prior work on VQC-based reinforcement learning in low-dimensional environments.}
  \label{tab:related}
  \footnotesize
  \begin{tabular}{lp{4.2cm}cc}
    \hline
    Work & Architecture & DQN? & Variants \\
    \hline
    Chen et al.~\cite{chen2020vqc} & 4--8q, 1--4L (CartPole, Acrobot) & \checkmark & No \\
    Jerbi et al.~\cite{jerbi2021pqc} & Policy-gradient VQC (CartPole) & No & Policy \\
    Lockwood \& Si~\cite{lockwood2021rl} & 4q, 2L (CartPole) & \checkmark & No \\
    Skolik et al.~\cite{skolik2022quantum} & 4q, re-upload (CartPole, Acrobot) & \checkmark & No \\
    \textbf{This work} & 4q, L1/L3 (CartPole) & \checkmark & \textbf{Double, Duel, Distill} \\
    \hline
  \end{tabular}
\end{table}

The VQC-based DQN architecture in this paper follows the angle-encoding plus
ring-CNOT design common to prior work, with the addition of batched QNode
evaluation for simulation efficiency and systematic evaluation of DQN variants.
The batched QNode optimization replaces per-sample circuit calls with a single
batched call using PennyLane parameter broadcasting, a practical contribution
relevant to any VQC-in-the-loop RL simulation.
